% GENERATED by scripts/render-cv.js — DO NOT EDIT BY HAND.
% Source: data/experience.json via cv/templates/biosketch-experience.tex.njk. Run `npm run render`.

{\bf AI Scientist -- Geospatial \& Remote Sensing, Woodwell Climate Research Center, MA, USA (remote)} \hfill {Jan 2022 - }

\begin{itemize}
    \item Project Lead: the pan-Arctic Retrogressive Thaw Slumps (RTS) mapping initiative, applying state-of-the-art deep learning techniques to satellite imagery for detecting and quantifying permafrost thaw features, directly contributing to climate feedback modelling and carbon cycle research
    \item Designed and built the ARTS scientific dataset, establishing a benchmark training resource for deep learning applications in environmental monitoring
    \item Published the first circumpolar retrogressive thaw slump probability map and polygon inventory (NSF Arctic Data Center, doi:10.18739/A26970107), produced by a 41.5-million-tile deep learning inference sweep of 2025 PlanetScope imagery: 19,068 quality-controlled slump polygons covering 529.7 km$^2$, released as cloud-optimised rasters and open vectors
    \item Developed innovative machine learning workflows for multimodal data integration, combining various satellite imagery types and field observations to enhance feature detection accuracy
    \item Developed A Zero-shot, interactive segmentation workflow for satellite image labelling using the Segment Anything Model (SAM)
    \item Collaborated with UCONN in using Location Embeddings to improve satellite image segmentation with vision foundation models.
    \item Collaborated with ASU in using big geospatial data to produce pan-Arctic statistics.
    \item Collaborated with UIUC in optimising the workflow of the Deep Learning model inference in big geospatial data (Arctic-scale) in GCP.
    \item Collaborated with ASU in modelling abrupt thaw initialisation with remote sensing and deep neural networks.
    \item Collaborated in two wildfire projects using deep learning for wildfire mapping on satellite imagery and using environmental variables to predict wildfire using machine learning
    \item Collaborated in carbon flux time series imputation using machine learning
    \item Led the machine learning workshop series for Woodwell researchers in 2022 and 2023, providing training in advanced AI techniques for environmental science applications including image processing and time-series forecasting. Designed progressive curriculum materials and adapted teaching methods for diverse learner backgrounds.
    \item Mentored interns and research assistants, created training materials for interns and research assistants. Mentee Works were presented at the AGU annual conference.
\end{itemize}

{\bf Data Science Fellow, Faculty.ai, London, UK} \hfill {Sep 2021 - Dec 2021}

\begin{itemize}
    \item Completed intensive fellowship in Machine Learning and Artificial Intelligence, focusing on developing practical AI solutions for complex data challenges. Include trainings and courses for professional data scientist and an industry placement.
\end{itemize}

{\bf PhD Researcher, International Centre for Carbonate Reservoirs, Edinburgh, UK} \hfill {Sep 2017 - Sep 2020}

\begin{itemize}
    \item Funded by PetroBras, petrophysics and digital rock experimental study.
    \item Conducted doctoral research on multiphase fluid flow in porous media, developing novel computational methods for analyzing complex flow dynamics, funded by Petróleo Brasileiro
    \item Parcipitated in petrophysical experiments at the Diamond Light Source UK, the Advanced Photon Source Argonne National Laboratory, US and the Swiss Light Source Paul Scherrer Institute, CH
\end{itemize}

